\documentclass[12pt,oneside]{scrreport} % KOMA: mejor tipografía, buen control de capítulos

% --- Idioma y tipografía ---
\usepackage{newpxtext} % classy typography
\usepackage{fontspec}        % fuentes del sistema (LuaLaTeX)
\usepackage{mathtools}       % antes de unicode-math (requerido)
\usepackage{amsthm}          % antes de unicode-math para que thmtools detecte el backend
\usepackage{unicode-math}    % fuentes matemáticas unicode; provee \symbf; reemplaza amssymb
\usepackage{microtype}       % mejor espaciado
\usepackage[spanish,es-noshorthands]{babel}
\usepackage{csquotes}        % cargar después de babel

% --- Página ---
\usepackage[a4paper,margin=2.6cm]{geometry}

% --- Matemáticas ---
% mathtools cargado arriba (antes de unicode-math); amssymb reemplazado por unicode-math
\usepackage{thmtools}        % estilos/refinamientos para teoremas
\numberwithin{equation}{section}

% Teoremas y amigos (numera por sección; todos comparten contador)
\declaretheoremstyle[
	headfont=\bfseries,
	bodyfont=\itshape,
	spaceabove=6pt, spacebelow=6pt, headpunct=.
]{thmstyle}
\declaretheorem[style=thmstyle,numberwithin=section,name=Teorema]{theorem}
\declaretheorem[style=remark,sibling=theorem,name=Informal]{informal}
\declaretheorem[style=thmstyle,sibling=theorem,name=Lema]{lemma}
\declaretheorem[style=thmstyle,sibling=theorem,name=Proposición]{proposition}
\declaretheorem[style=definition,sibling=theorem,name=Definición]{definition}
\declaretheorem[style=remark,sibling=theorem,name=Ejemplo]{example}
\declaretheorem[style=remark,sibling=theorem,numbered=no,name=Ejercicio]{exercise}
\declaretheorem[style=remark,sibling=theorem,numbered=no,name=Insight]{insight}

\newenvironment{solution}{\begin{proof}[Solución]\let\qed\relax}{\end{proof}}

% --- Cajas de teoremas ---
\usepackage[skins,breakable]{tcolorbox}

\tcbset{
	thmbase/.style={
			enhanced,
			breakable,
			sharp corners,
			frame hidden,
			left=6pt, right=4pt, top=3pt, bottom=3pt,
			before skip=8pt, after skip=8pt,
		}
}

% % Azul — resultados (teoremas, lemas, proposiciones)
\tcolorboxenvironment{theorem}{thmbase,
	borderline west={3pt}{0pt}{blue!70!black}, colback=blue!5!white}
\tcolorboxenvironment{lemma}{thmbase,
	borderline west={3pt}{0pt}{blue!70!black}, colback=blue!5!white}
\tcolorboxenvironment{proposition}{thmbase,
	borderline west={3pt}{0pt}{blue!70!black}, colback=blue!5!white}

% Verde azulado — definiciones
\tcolorboxenvironment{definition}{thmbase,
	borderline west={3pt}{0pt}{teal!80!black}, colback=teal!5!white}

% Naranja — ejemplos
\tcolorboxenvironment{example}{thmbase,
	borderline west={3pt}{0pt}{orange!80!black}, colback=white}

% Gris — informal, insight
\tcolorboxenvironment{informal}{thmbase,
	borderline west={3pt}{0pt}{gray!55!black}, colback=gray!5!white}
\tcolorboxenvironment{insight}{thmbase,
	borderline west={3pt}{0pt}{gray!55!black}, colback=white}

% --- Referencias y PDF ---
\usepackage{hyperref}
\hypersetup{colorlinks=true, linkcolor=blue!60!black, urlcolor=blue!60!black}
\usepackage[capitalise,noabbrev,nameinlink]{cleveref}
% Spanish names for theorem-like environments:
\crefname{theorem}{Teorema}{Teoremas}
\crefname{lemma}{Lema}{Lemas}
\crefname{proposition}{Proposición}{Proposiciones}
\crefname{corollary}{Corolario}{Corolarios}
\crefname{definition}{Definición}{Definiciones}
\crefname{example}{Ejemplo}{Ejemplos}
% spanish Prueba instead of proof 
\renewcommand{\proofname}{Prueba}

% to reference steps in proofs/solutions
\newcounter{proofstep}
\newcommand{\step}[1]{%
	\refstepcounter{proofstep}\label{#1}\textbf{(\theproofstep) }\ignorespaces}

\AtBeginEnvironment{proof}{\setcounter{proofstep}{0}}
\AtBeginEnvironment{solution}{\setcounter{proofstep}{0}}

% Make \cref{<proofstep-label>} print "(n)"
\crefformat{proofstep}{(#2#1#3)}

% Let TeX fix small overruns instead of whining
%\emergencystretch=2em   % extra stretch on the last pass (very effective)
%\hfuzz=2pt              % ignore overfull < 2pt (tune to 3pt if you like)
%\hbadness=3000          % reduce underfull warnings sin silenciarlos todos

% to use enum with letters
\usepackage{enumitem}
% --- Gráficos ---
\usepackage{graphicx}
\usepackage{subcaption}
\graphicspath{{figures/}}
% fancy plots
% \usepackage{tikz}
% \usetikzlibrary{arrows.meta,calc}
% --- Macros del curso ---
\newcommand{\N}{\mathbb{N}}
\newcommand{\Z}{\mathbb{Z}}
\newcommand{\Q}{\mathbb{Q}}
\newcommand{\R}{\mathbb{R}}
\newcommand{\C}{\mathbb{C}}
\newcommand{\K}{\mathbb{K}}
\DeclareMathOperator{\im}{Im}
\DeclareMathOperator{\tr}{trace}
\DeclareMathOperator{\rk}{rank}
\newcommand{\spn}[1]{\operatorname{span}\left(#1\right)}
\DeclarePairedDelimiter{\norm}{\lVert}{\rVert}
\DeclarePairedDelimiter{\ip}{\langle}{\rangle}
\DeclarePairedDelimiter{\abs}{|}{|}
\DeclarePairedDelimiter{\parens}{(}{)}
\newcommand{\fullref}[1]{\cref{#1} (\nameref{#1})}
\newcommand{\Lin}[1]{\mathcal{L}(#1)}
\newcommand{\adj}[1]{#1^{\ast}}

% tablas
\usepackage{booktabs}
\usepackage{tabularx}
